% Seção 3.6 -- Investigação adicional (notebook 6_investigacao_picos_sp.ipynb)
% Pacotes usados: graphicx, booktabs, amsmath (e babel em português).
% Figuras: subir para o Overleaf na pasta figuras/ os arquivos
%   results/figures/picos_sp_registros.png
%   results/figures/picos_sp_histograma.png

\section{Investigação adicional: velocidade de operação e picos espectrais}
\label{sec:investigacao}

\subsection*{Pergunta e justificativa}

\textbf{Como varia, entre os níveis de \texttt{SP}, a frequência do pico principal de 30 a 200\,Hz, e essa variação é a mesma em hélices saudáveis e defeituosas?}
Nos espectros da Seção~3.3, o maior pico de baixa frequência pareceu subir com o \texttt{SP} e ficar um pouco abaixo na hélice defeituosa, mas isso foi visto só em três pares ilustrativos.
Além disso, as características de amplitude (RMS, pico) mudam de sentido entre os níveis de \texttt{SP} (Seção~3.4).
Uma frequência poderia ser uma característica mais estável.
A documentação do DronePropB não informa a rotação correspondente a cada \texttt{SP} nem a origem física dos picos, então os picos são tratados apenas como componentes do espectro, sem atribuição de causa.

\subsection*{Método}

Para cada uma das 2370 janelas de 2\,s ($\Delta f = 0{,}5$\,Hz, janela de Hann, sem a média), calculou-se $f_1$, a frequência de maior magnitude entre 30 e 200\,Hz.
Em \texttt{SP1} e \texttt{SP2}, $f_1$ salta entre um pico e outro com o dobro da frequência (Figura~\ref{fig:picos_hist}, linha de cima): $\approx 53$ e $\approx 106$\,Hz em \texttt{SP1}, e $\approx 95$ e $\approx 190$\,Hz em \texttt{SP2}.
Por isso, definiu-se um pico de base $f_\text{base}$: se existe um pico em $f_1/2 \pm 2$\,Hz com pelo menos 30\% da magnitude de $f_1$, então $f_\text{base}$ é esse pico; caso contrário, $f_\text{base} = f_1$.
Cada registro é resumido pela mediana de suas janelas, porque as janelas de um mesmo registro não são independentes.
A separação entre condições é medida pela AUC (estatística $U$ de Mann--Whitney dividida por $n_\text{def}\,n_\text{saud}$) em cada \texttt{SP}, com 6 saudáveis e 27 defeituosos por nível.
AUC $< 0{,}5$ indica frequência menor na defeituosa.

\begin{figure}[htbp]
  \centering
  \includegraphics[width=\linewidth]{figuras/picos_sp_histograma.png}
  \caption{Distribuição por janela de $f_1$ (em cima) e de $f_\text{base}$ (embaixo), em \% das janelas de cada condição. Em \texttt{SP1} e \texttt{SP2}, $f_1$ forma dois grupos com uma frequência o dobro da outra. A correção leva a maior parte das janelas ao grupo mais baixo.}
  \label{fig:picos_hist}
\end{figure}

\subsection*{Resultados}

A Tabela~\ref{tab:picos} e a Figura~\ref{fig:picos_reg} resumem $f_\text{base}$ por registro.
O pico de base sobe com o \texttt{SP} nas duas condições, e os picos em 2 e 3 vezes $f_\text{base}$ acompanham o deslocamento.
A ordem $\texttt{SP1} < \texttt{SP2} < \texttt{SP3}$ vale para 97 dos 99 registros.
As duas exceções são saudáveis (\texttt{F0\_SV0\_SP2\_C3\_R2} e \texttt{F0\_SV0\_SP3\_C3\_R2}) com metade das janelas em cada um de dois picos, em que a mediana cai entre eles.
O crescimento não é proporcional ao índice do \texttt{SP}: em relação a \texttt{SP1}, as razões são $\approx 1{,}8$ e $\approx 2{,}4$.

\begin{table}[htbp]
  \centering
  \caption{$f_\text{base}$ por \texttt{SP} e condição (mediana dos registros). AUC $< 0{,}5$: menor na defeituosa. A última coluna usa só os saudáveis gravados em setembro, o mesmo mês dos defeituosos.}
  \label{tab:picos}
  \begin{tabular}{cccccc}
    \toprule
    \texttt{SP} & Saudável (Hz) & Defeituosa (Hz) & Diferença & AUC & AUC (só setembro) \\
    \midrule
    1 & 56,2  & 53,0  & $-5{,}8\%$ & 0,22 & 0,20 \\
    2 & 100,0 & 95,5  & $-4{,}5\%$ & 0,15 & 0,19 \\
    3 & 133,5 & 130,2 & $-2{,}4\%$ & 0,38 & 0,21 \\
    \bottomrule
  \end{tabular}
\end{table}

\begin{figure}[htbp]
  \centering
  \includegraphics[width=\linewidth]{figuras/picos_sp_registros.png}
  \caption{$f_\text{base}$ de cada registro (mediana das janelas), por \texttt{SP}. Os saudáveis estão separados por mês de gravação, e os defeituosos são todos de setembro. Traço preto: mediana do grupo.}
  \label{fig:picos_reg}
\end{figure}

\textbf{A hélice defeituosa tem $f_\text{base}$ de 2 a 6\% menor nos três níveis de \texttt{SP}, e varia mais entre registros.}
Em \texttt{SP1}, os seis saudáveis ficam entre 56,0 e 56,8\,Hz, enquanto os defeituosos vão de 47 a 60\,Hz, e 74\% deles ficam abaixo do menor saudável.
Ao contrário do RMS, o sentido da diferença é o mesmo nos três níveis.

\textbf{A data de gravação não parece explicar a diferença.}
Os saudáveis gravados em setembro têm $f_\text{base}$ praticamente igual aos de novembro em \texttt{SP1} e \texttt{SP2} (56,4 contra 56,2\,Hz e 99,8 contra 100,0\,Hz), e usá-los sozinhos mantém o sentido da AUC (Tabela~\ref{tab:picos}).
Isso contrasta com a média do sinal, que acompanha a data (Seção~3.4).
Com dois saudáveis de setembro por \texttt{SP}, essa é uma verificação parcial.

\textbf{O resultado não depende do estimador.}
Com $f_1$ sem correção e com limiares de 20\% e 50\%, a defeituosa continua abaixo da saudável nos três níveis.
Sem correção, a AUC fica mais próxima de 0,5 em \texttt{SP1} e \texttt{SP2} (0,41 e 0,27), porque os saltos para o dobro da frequência misturam os registros.

\subsection*{Limitações}

Não se sabe a rotação real de cada \texttt{SP}, nem por que a defeituosa tem pico de base menor.
$f_\text{base}$ é o pico mais baixo de uma série dentro de 30--200\,Hz, e não necessariamente a fundamental.
A correção por $f_1/2$ é heurística.
Há poucos saudáveis por \texttt{SP} (6), as repetições do mesmo grupo são parecidas entre si, e a classe defeituosa mistura tipos (\texttt{F1}--\texttt{F3}) e severidades (\texttt{SV1}--\texttt{SV3}) não documentados no DronePropB.
As AUCs são, portanto, descritivas.
Para a classificação (Seção~3.5), $f_\text{base}$ é uma característica espectral mais estável que $f_1$.
